\section{The order of choices changes the meaning}\label{ch02:sec:quantifiers}
A \emph{universal} statement, such as ``every integer divisible by four is even,'' claims that every allowed input has a property. A proof cannot check infinitely many integers one at a time; instead, as in the proof of Proposition~\ref{thm:odd-sum}, it works with an \emph{arbitrary} input, one about which we know nothing except that it is allowed, so the same argument applies to each input. An \emph{existence} statement, such as ``some integer is both even and divisible by three,'' claims that at least one allowed input has the property. A single example, such as $6$, proves it; there may be others, such as $12$.

The symbol $\forall$ means ``for every,'' and the symbol $\exists$ means ``there exists'' or ``there is at least one.'' They are called \emph{quantifiers}, because they say for how many values of a letter a condition must hold. We write $x\in\RR$, read ``$x$ in $\RR$,'' to say that $x$ belongs to the set $\RR$ of real numbers. The set that a quantified letter ranges over is called its \emph{domain}; in ``$\forall x\in\RR$'' the domain of $x$ is $\RR$. When several quantifiers appear in one statement, their order matters. Compare
\begin{align*}
&\forall x\in\RR\ \exists y\in\RR:\quad y>x,\tag{A}\\
&\exists y\in\RR\ \forall x\in\RR:\quad y>x.\tag{B}
\end{align*}
Read aloud, (A) says ``for every real number $x$, there is a real number $y$ such that $y>x$,'' and (B) says ``there is a real number $y$ such that, for every real number $x$, $y>x$.'' The colon separates the quantifiers from the condition that the chosen values must satisfy.

It helps to read a quantified statement as a game between a skeptic, who doubts the statement, and a prover, who defends it. The quantifiers are read from left to right, and each one is a move: at a $\forall$, the skeptic picks any member of the domain, and at an $\exists$, the prover picks a value. The prover wins if the condition after the colon holds for the values picked, and the statement is true exactly when the prover has a way to win however the skeptic plays.

Statement (A) is true. The skeptic picks $x$ first, and the prover may then respond with a $y$ that suits it: $y=x+1$ works. The response depends on the input: for $x=3$ the prover answers $4$, and for $x=-10$ it answers $-9$. Statement (B) is false. The prover must commit to a single $y$ \emph{first}, and whatever $y$ is chosen, the skeptic's reply $x=y+1$ defeats it, because $y>y+1$ is false.

In general, a value chosen at an $\exists$ may depend on every value chosen before it, but not on anything chosen after it; this is the dependence of Section~\ref{ch01:sec:dependence}, where the witnesses $k$ and $\ell$ were determined by $a$ and $b$. Swapping two neighboring quantifiers of the same kind, as in $\forall x\in\RR\ \forall y\in\RR$, does not change the meaning, as long as the domain of neither letter refers to the other. (In ``for every integer $N$ and every integer $n\ge N$,'' the domain of $n$ depends on $N$, so $N$ must come first.) When a proof that chooses several quantities becomes confusing, ask: \emph{at what stage is this quantity chosen, and which other quantities may it depend on?}

\subsection*{Tolerances and the inequality rules they need}
In later chapters we often need one number to be within some allowed error of another. A \emph{tolerance} is a positive amount of error that we allow, traditionally written $\eps$, the Greek letter epsilon. A statement that begins ``for every positive $\eps$,'' written $\forall\eps>0$, must hold for $\eps=1/10$, for $\eps=1/1000$, and for every other positive real number.

Tolerances lead to inequalities, where the sign of a number matters. Three rules cover almost everything we need.
\begin{itemize}
\item Adding the same number to both sides keeps the direction of an inequality: from $2<5$ we get $2+7<5+7$.
\item Multiplying or dividing both sides by the same \emph{positive} number keeps the direction: multiplying $2<5$ by $3$ gives $6<15$.
\item Multiplying or dividing both sides by the same \emph{negative} number reverses the direction: multiplying $2<5$ by $-1$ gives $-2>-5$.
\end{itemize}
The same three rules hold with $\le$ in place of $<$.

The \emph{reciprocal} of a nonzero number $a$ is $1/a$. Suppose $0<a<b$. The product $ab$ is positive, so by the second rule we may divide both sides of $a<b$ by $ab$:
\[
\frac{a}{ab}<\frac{b}{ab},\qquad\text{that is,}\qquad \frac1b<\frac1a.
\]
So for positive numbers, taking reciprocals reverses the order: from $2<5$ we get $1/5<1/2$. The same argument shows that $0<a\le b$ implies $1/b\le1/a$. The rule can fail when the numbers are not both positive: $-1<2$, but $1/(-1)=-1$ is smaller than $1/2$.

Given a tolerance $\eps$, we will often want a positive integer $N$ with $1/N<\eps$, and for this it is enough that $N>1/\eps$. Such an integer always exists, by the following property of the real numbers, which we accept without proof; it is the second of the starting facts listed in \emph{To the reader}.
\begin{center}
\emph{The Archimedean property: for every real number $r$, there is an integer $N$ with $N>r$.}
\end{center}

\begin{example}\label{ch02:ex:archimedes}
We prove that for every positive real number $\eps$, there is a positive integer $N$ with $1/N<\eps$. Let $\eps$ be a positive real number. By the Archimedean property, there is an integer $N$ with $N>1/\eps$. This $N$ is positive, since $N>1/\eps>0$ (divide both sides of $1>0$ by $\eps$ to get $1/\eps>0$). As $1/\eps$ and $N$ are both positive, taking reciprocals reverses the inequality $1/\eps<N$:
\[
\frac1N<\frac{1}{1/\eps}=\eps.
\]
The order of the quantifiers allows $N$ to depend on $\eps$, and it has to: a smaller tolerance may need a larger $N$. For instance, if $\eps=0.003$, then $1/\eps=333.33\ldots$, so $N=334$ works, with $1/334=0.002994\ldots$, while $N=333$ fails, because $1/333=0.003003\ldots$ is larger than $0.003$.
\end{example}

\pauseq{Can a single positive integer $N$ satisfy $1/N<\eps$ for every positive $\eps$?}{No. Once $N$ is fixed, the skeptic may choose the positive number $\eps=1/(2N)$. Then $1/N=2\eps>\eps$, so $1/N<\eps$ fails. This claim has the shape of (B): one $N$ chosen before $\eps$. The true statement in Example~\ref{ch02:ex:archimedes} has the shape of (A): $N$ is chosen after $\eps$ is known.}
